\documentclass[12pt]{article}
\usepackage{latexblog}

% ============================================================
% Blog Metadata
% ============================================================
\blogtitle{理解Java并发(2)：ThreadPoolExecutor解析}
\blogdate{2019-10-31}
\blogtags{刨根问底, ~理解Java并发, 并发编程, Concurrent}
\bloglang{zh}
\blogsummary{}

\begin{document}
\makeblogtitle

使用多线程技术可以有效的利用CPU时间，在同一个时间内完成更多的任务，但同时值得注意的是，线程创建本身也是有开销的，线程池使得我们可以重复的利用已经存在的线程，从而节省这一部分的开销，提高程序的效率。

\section{线程数的限制}\label{ux7ebfux7a0bux6570ux7684ux9650ux5236}

首先一个问题是，我们在创建新的线程的时候，是不是线程越多就越好呢？实际上是不可能无限的创建新的线程的，总会有个限制，那么问题是这个限制是多大，或者说取决于什么呢？

\subsection{操作系统的最大线程数}\label{ux64cdux4f5cux7cfbux7edfux7684ux6700ux5927ux7ebfux7a0bux6570}

首先希望搞清楚的一个问题就是，到底我们可以创建多少个线程呢？在linux上，可以通过以下的方式查看系统的最大线程数限制：

\begin{verbatim}
# cat /proc/sys/kernel/threads-max
15734
# ulimit -v
unlimited
\end{verbatim}

据说是按照这个公式计算出来的:

\begin{Shaded}
\begin{Highlighting}[]
\NormalTok{max\_threads }\OperatorTok{=}\NormalTok{ mempages }\OperatorTok{/} \OperatorTok{(}\DecValTok{8} \OperatorTok{*}\NormalTok{ THREAD\_SIZE }\OperatorTok{/}\NormalTok{ PAGE\_SIZE}\OperatorTok{);}
\end{Highlighting}
\end{Shaded}

在windows上也比较类似，总结来说就是，每个系统的最大线程数都不尽相同，不仅与系统有关还与内存大小以及用户的设置有关系。

\subsection{JVM限制}\label{jvmux9650ux5236}

JVM本身貌似没有对线程数进行限制，但同样不能无限制的创建线程否则会出现\texttt{java.lang.OutOfMemoryError:\ unable\ to\ create\ new\ native\ thread}。在JVM中有以下的一些参数可能会影响能创建的线程数：

\begin{itemize}
\tightlist
\item
  -Xms 设置堆的最小值
\item
  -Xmx 设置堆的最大值
\item
  -Xss 设置每个线程的stack大小
\end{itemize}

因为一个机器上的内存是一定的，所以如果\texttt{-Xss}设置的越大，单个线程所占用的栈空间越大，那么能创建的线程数就越少。一个比较有趣的事实是，能创建的最大线程数是跟\texttt{-Xmx}的值负相关的，即你设置的堆越大，反而能创建的最大线程数越少！这是别人的测试结果：

\begin{verbatim}
2 mb --> 5744 threads
4 mb --> 5743 threads
...
768 mb --> 3388 threads
1024 mb --> 2583 threads
\end{verbatim}

原因就是堆空间越大，那么机器上剩下的内存就越少，即可以用来分配给线程栈上的内存就越少，所以会出现这样的结果。

\section{线程池}\label{ux7ebfux7a0bux6c60}

在Java中线程的启动和停止是有开销的。这个开销主要包括：

\begin{itemize}
\tightlist
\item
  为线程开辟栈空间（例如OpenJDK6在Linux上会使用\texttt{pthread\_create}来创建线程，内部使用\texttt{mmap}分配内存)
\item
  通过操作系统的调用来创建和注册本地线程
\item
  保存线程的相关信息（JVM/native thread descriptors)到JVM中
\end{itemize}

根据网上的测试来看，通常使用线程池可以获得大幅的性能提升（亲测至少15倍）。而使用线程池相当于重用了已有的线程，避免了这部分开销。当任务越多越频繁的情况下，这部分开销越不可小觑。

\subsection{ThreadPoolExecutor}\label{threadpoolexecutor}

ThreadPoolExecutor 是一个利用线程池技术实现的多任务处理器，它的申明如下：

\begin{Shaded}
\begin{Highlighting}[]
\KeywordTok{public} \BuiltInTok{ThreadPoolExecutor}\OperatorTok{(}\DataTypeTok{int}\NormalTok{ corePoolSize}\OperatorTok{,}
                            \DataTypeTok{int}\NormalTok{ maximumPoolSize}\OperatorTok{,}
                            \DataTypeTok{long}\NormalTok{ keepAliveTime}\OperatorTok{,}
                            \BuiltInTok{TimeUnit}\NormalTok{ unit}\OperatorTok{,}
                            \BuiltInTok{BlockingQueue}\OperatorTok{\textless{}}\BuiltInTok{Runnable}\OperatorTok{\textgreater{}}\NormalTok{ workQueue}\OperatorTok{)} \OperatorTok{\{}
    \KeywordTok{this}\OperatorTok{(}\NormalTok{corePoolSize}\OperatorTok{,}\NormalTok{ maximumPoolSize}\OperatorTok{,}\NormalTok{ keepAliveTime}\OperatorTok{,}\NormalTok{ unit}\OperatorTok{,}\NormalTok{ workQueue}\OperatorTok{,}
            \BuiltInTok{Executors}\OperatorTok{.}\FunctionTok{defaultThreadFactory}\OperatorTok{(),}\NormalTok{ defaultHandler}\OperatorTok{);}
\OperatorTok{\}}
\end{Highlighting}
\end{Shaded}

乍一看有很多个参数，那么该如何去配置呢？

\subsubsection{CorePoolSize / MaximumPoolSize}\label{corepoolsize-maximumpoolsize}

线程池会根据这两个参数去管理池中的线程。当一个新的任务提交的时候，会遵循如下的规则：

\begin{itemize}
\tightlist
\item
  如果池中线程数小于corePoolSize，哪怕有空闲的线程也会创建一个新的线程来
\item
  当吃中线程数超过corePoolSize但是小于MaximumPoolSize的时候，只有当workQueue满的时候才会创建新的线程
\end{itemize}

所以当corePoolSize和MaximumPoolSize一样的时候，实际上就是一个固定大小的线程池，相当于使用\texttt{Executors.newFixedThreadPool}:

\begin{Shaded}
\begin{Highlighting}[]
\KeywordTok{public} \DataTypeTok{static} \BuiltInTok{ExecutorService} \FunctionTok{newFixedThreadPool}\OperatorTok{(}\DataTypeTok{int}\NormalTok{ nThreads}\OperatorTok{)} \OperatorTok{\{}
    \ControlFlowTok{return} \KeywordTok{new} \BuiltInTok{ThreadPoolExecutor}\OperatorTok{(}\NormalTok{nThreads}\OperatorTok{,}\NormalTok{ nThreads}\OperatorTok{,}
                                  \DecValTok{0L}\OperatorTok{,} \BuiltInTok{TimeUnit}\OperatorTok{.}\FunctionTok{MILLISECONDS}\OperatorTok{,}
                                  \KeywordTok{new} \BuiltInTok{LinkedBlockingQueue}\OperatorTok{\textless{}}\BuiltInTok{Runnable}\OperatorTok{\textgreater{}());}
\OperatorTok{\}}
\end{Highlighting}
\end{Shaded}

池中的线程默认只要当提交了新任务的时候才会创建，如果希望提前创建线程可以使用\texttt{prestartCoreThread}或者\texttt{prestartAllCoreThreads}。

\subsubsection{Keep-alive 时间}\label{keep-alive-ux65f6ux95f4}

当池中的线程数多余corePoolSize的时候，超出部分的线程会在空闲一段时间之后被终止掉，这个时间就是keepAliveTime。如果设置为0那么一旦超出部分运行结束之后就会被终止掉，反之如果设置为\texttt{Long.MAX\_VALUE}那么空闲线程就会一直存活。

默认情况下，只有超出corePoolSize的线程才会受到这个存活时间的影响，如果希望对于核心线程也能超时终止，那么可以使用\texttt{allowCoreThreadTimeOut}来控制。

\subsubsection{workQueue}\label{workqueue}

工作队列用来持有提交的任务。规则如下：

\begin{itemize}
\tightlist
\item
  如果当前池中的线程少于corePoolSize，则创建新的线程
\item
  如果大于corePoolSize，则倾向于将任务加入到workQueue中
\item
  如果无法将任务加入到队列中，则会创建新的线程，直到池中的线程数达到maximumPoolSize
\item
  如果超过maximumPoolSize，那么将会拒绝提交的任务
\end{itemize}

对于队列的选择也可以使用不同的策略：

\begin{itemize}
\tightlist
\item
  使用\texttt{SynchronousQueue}可以直接将任务从队列转手到线程池，这个参数要配合将maximumPoolSize设置为无限大来配合使用。因为这个朝这个队列中插入一条数据将会阻塞一直到它被消费，也就是说读写操作要配套，实际上就是进行了一个数据交换，根本没有在队列中实际存储任务。如果maximumPoolSize太小可能会导致任务被拒绝。
\item
  使用无界的队列例如\texttt{LinkedBlockingQueue}，那么一旦线程超过corePoolSize的时候新线程都会被加入到队列中，也就是说maximumPoolSize根本不会生效了。
\item
  使用有界队列例如\texttt{ArrayBlockingQueue}，超过队列数的新任务将创建新的线程。那么这时候队列大小和线程数上限需要权衡配合。
\end{itemize}

References:

\begin{itemize}
\tightlist
\item
  \href{https://stackoverflow.com/questions/344203/maximum-number-of-threads-per-process-in-linux}{Maximum number of threads per process in Linux?}
\item
  \href{https://eknowledger.wordpress.com/2012/05/01/max-number-of-threads-per-windows-process/}{Max Number of Threads Per Windows Process}
\item
  \href{http://baddotrobot.com/blog/2009/02/26/less-is-more/}{Less is More}
\item
  \href{https://stackoverflow.com/questions/5483047/why-is-creating-a-thread-said-to-be-expensive}{Why is creating a Thread said to be expensive?}
\end{itemize}


\end{document}
